\begin{proposition}[The sigma expression and third-kind period]\hypertarget{tag-0017}{}\label{ari:C.7.1}\hypertarget{tag-IA01}{}\label{ari:C.7.1:lift-transport}
\StatementRecord{ari:C.7.1}{proved}{0017}{\SeriesRef{ari:eq:cp-periods}; \citet[Definitions 2.3--2.4 and
Section 6]{bertolin2022}}

Let $P\in E(\overline\Q)$, with $P\ne O$ and $2P\ne O$, and
let $u$ be an elliptic logarithm of $P$. Then
\begin{equation}\hypertarget{tag-003K}{}\label{ari:eq:cp-sigma}
 2\log N(P)=\frac16\log|\Delta|
       +\Re\{\Lambda_1+u(\zeta(u)-\eta(u))\},
 \qquad \Lambda_1=2\log\sigma(u)-u\zeta(u).
\end{equation}
For a fixed elliptic logarithm $u$, changing the branch of
$\log\sigma(u)$ changes the expression inside the real part by an
integer multiple of $4\pi i$. Put
$B(u)=\Lambda_1+u(\zeta(u)-\eta(u))$. Under a lattice translation,
\[
 B(u+\omega)-B(u)\equiv
 u\eta(\omega)-\omega\eta(u)+2\log\chi(\omega)
 \pmod{4\pi i\Z},\qquad \omega\in\Lambda,
\]
where $\chi(\omega)\in\{1,-1\}$ is the sigma semicharacter.
This difference is purely imaginary, so the displayed real part
is independent of both choices.
There is an extension $G_P$ of $E$ by $\mathbb G_m$, defined
over $\overline\Q$ and corresponding, under the principal
polarization, to $P$, for which $\Lambda_1$, modulo a logarithm
of an algebraic number, is the negative of a relative period
of $[\Z\to G_P]$. Changing the third-kind differential while retaining its simple
poles and residues adds an algebraic multiple of $u$.
Adding an exact algebraic differential to a relative-integral
representative adds algebraic endpoint differences when both
endpoints are regular.
\end{proposition}


\begin{proof}
\begin{enumerate}[label=\arabic*.,leftmargin=*]
\item The sigma--theta product, or its divisor and automorphy law
together with its derivative at zero, gives
\begin{equation}\hypertarget{tag-003F}{}\label{ari:eq:cp-klein}
 N(P)=|\Delta|^{1/12}|\sigma(u)|
                   \exp\!\left(-\frac12\Re(u\eta(u))\right).
\end{equation}
To check the constant, for $\Lambda=\Z+\tau\Z$ one uses
$\sigma(z)=e^{\eta_1z^2/2}\vartheta_1(z,\tau)/(2\pi\eta(\tau)^3)$,
$\Delta=(2\pi)^{12}\eta(\tau)^{24}$, and
$\eta(z)=\eta_1z-2\pi i\Im z/T$. Substitution gives exactly
the defining Gaussian in $N$. Homogeneity gives the general
lattice formula. Taking logarithms proves \eqref{ari:eq:cp-sigma}.
The automorphy law
$\sigma(u+\omega)=\chi(\omega)
 e^{\eta(\omega)(u+\omega/2)}\sigma(u)$ and real-linearity of
$\eta$ give the stated lift transport. Explicitly, writing
$u=a\omega_1+b\omega_2$ with $a,b\in\R$ and
$\omega=m\omega_1+n\omega_2$ with $m,n\in\Z$, Legendre's
relation gives
$u\eta(\omega)-\omega\eta(u)=2\pi i(bm-an)$.
This also proves that the transported difference is purely imaginary.

\item We make the period identification explicit. Put
\[
 f_u(z)=\frac{\sigma(z+u)}{\sigma(z)\sigma(u)}e^{-\zeta(u)z}.
\]
Its logarithmic differential is
\[
 \frac{df_u}{f_u}
 =\frac12\frac{\wp'(z)-\wp'(u)}{\wp(z)-\wp(u)}\,dz.
\]
This is the pullback of an algebraic differential with residues
$+1$ at $-P$ and $-1$ at $O$.
Two algebraic differentials with precisely these simple poles
and residues differ by a holomorphic differential, hence by
$c\,dz$ with $c\in\overline{\mathbb Q}$.
Integration gives $cu$; integration of $df$ gives the stated
endpoint difference.
Here the principal polarization identifies $P$ with
$t_P^*\mathcal O_E(O)\otimes\mathcal O_E(O)^{-1}
=\mathcal O_E((-P)-(O))$; this fixes the sign of the extension class. The apparent $0/0$ at $z=u$
has a removable singularity when $P$ is not $2$-torsion.
The rational cocycle $f_u(z+z')/(f_u(z)f_u(z'))$ defines
the algebraic extension; in its rational trivialization the
exponential map has the form
\[
 \exp_{G_P}(z,v)=(\exp_E z,e^v f_u(z)).
\]
Its periods include
\begin{equation}\hypertarget{tag-003G}{}\label{ari:eq:cp-periods}
 (0,2\pi i),\qquad
 (\omega_j,\zeta(u)\omega_j-\eta_j u),\quad j=1,2.
\end{equation}
These signs follow directly from the sigma automorphy law.
The identity
$\sigma(2u)=-\wp'(u)\sigma(u)^4$ implies
\[
 \log f_u(u)=\Lambda_1+\log(-\wp'(u)).
\]
For an algebraic lift $R=(P,r)$, $r\in\overline\Q^{\times}$,
the second coordinate of $\log_{G_P}R$ is
\begin{equation}\hypertarget{tag-003L}{}\label{ari:eq:cp-v}
 v=\log r-\Lambda_1-\log(-\wp'(u)).
\end{equation}
It is a period of the relative homology defining the indicated
$1$-motive. This also avoids treating an integral with an endpoint
at a residue pole as an unregularized ordinary integral.
The description agrees with the generalized logarithm and
Serre function in \citet[Section~2.3, equations~(2.10)--(2.13), pp.~7--8;
Definitions~2.3--2.4, p.~8; Section~6.1, pp.~26--27,
arXiv v2]{bertolin2022}; the signs in
\eqref{ari:eq:cp-periods} use the displayed exponential convention.
\end{enumerate}
\end{proof}